\section{实践建议：什么时候用哪一种计划}

结合本讲内容，可以把两种机制分工如下：

\begin{table}[H]
\centering
\begin{tabular}{lll}
\toprule
机制 & 主要用途 & 典型场景 \\
\midrule
Plan Mode & 澄清需求、形成规格 & 用户目标模糊、任务路径多、需要选择取舍 \\
\texttt{update\_plan} & 执行中进度管理 & 已经开始做任务，需要展示状态和动态调整 \\
文件化 todo & 跨会话持久化 & 长任务、批量任务、需要恢复和审计 \\
\bottomrule
\end{tabular}
\caption{三种计划载体的分工。}
\end{table}

对于 coding agent 来说，最稳的工作流通常是：先探索代码库，必要时进入 Plan Mode 澄清规格；执行阶段用 \texttt{update\_plan} 管理当前进度；如果任务跨度很大，则额外维护文件化计划。

\begin{warningbox}{计划不是越细越好}
计划的价值在于减少不确定性，而不是制造仪式感。过细的计划会让 agent 在低价值步骤上浪费上下文；过粗的计划又无法指导执行。好的计划应该只捕捉真正影响执行顺序和风险控制的步骤。
\end{warningbox}

\subsection{本章小结}

Plan Mode、\texttt{update\_plan} 和文件化 todo 是三个层次。它们分别解决协作澄清、执行状态和长期持久化问题。

